\section{Despliegue y operación}

\begin{figure}[htbp]
\centering
\resizebox{0.96\textwidth}{!}{% Figura: despliegue real (desde 2026-08-25). Los servicios estan en loopback y lo
% unico publicado es nginx con TLS. VisioneFlow consume la API como un tercero mas.
\begin{tikzpicture}[node distance=7mm]

% --- Quien consume, desde internet ------------------------------------------
\node[compx,text width=34mm] (nav) {Navegador\\\nota{persona del equipo}};
\node[compx,text width=34mm,below=14mm of nav] (vf) {VisioneFlow\\\nota{nodo de agente}\\\nota{y disparo cada hora}};
\node[grupo,fit=(nav)(vf),label={[titgrupo]above:INTERNET}] (gnet) {};

% --- El servidor ------------------------------------------------------------
\node[comp,text width=30mm,right=26mm of nav,yshift=-9mm] (ngx)
  {nginx\\\nota{TLS, puertos 80 y 443}\\\nota{limita el acceso}};

\node[comp,font=\scriptsize,text width=44mm,right=30mm of ngx,yshift=16mm] (con)
  {Consola\\\nota{\cd{127.0.0.1:3001}}\\\nota{contraseña y cookie firmada}};
\node[comp,font=\scriptsize,text width=44mm,below=of con] (pre)
  {\cd{predictivo-predictivo-1}\\\nota{\cd{127.0.0.1:8000}}};
\node[comp,font=\scriptsize,text width=44mm,below=of pre] (his)
  {\cd{historico-historico-1}\\\nota{\cd{127.0.0.1:8010}}};
\node[agro,font=\scriptsize,text width=44mm,below=of his] (pg)
  {\cd{agrodash-pg}\\\nota{réplica del volcado}\\\nota{\cd{127.0.0.1:5433}}};

\node[comp,font=\scriptsize,text width=32mm,below=14mm of ngx] (tim)
  {temporizadores systemd\\\nota{\cd{predictivo-etl} 15 min}\\\nota{\cd{predictivo-refresh} 6 h}};

\node[grupo,fit=(ngx)(con)(pre)(his)(pg)(tim),
      label={[titgrupo]above:EC2 VisioneMetrics, apagada de 19:00 a 07:00}] (gsrv) {};

% --- Almacenamiento ---------------------------------------------------------
\node[almacen,font=\scriptsize,text width=34mm,right=14mm of his] (sb)
  {Supabase\\\nota{RLS activo}};

% --- Flujos -----------------------------------------------------------------
\draw[flecha] (nav.east) -- node[etiq,above,sloped]{HTTPS} (ngx.west);
\draw[flecha] (vf.east)  -- node[etiq,below,sloped]{HTTPS + clave} (ngx.west);

\draw[flecha] (ngx.east) -- node[etiq,pos=0.62,above,sloped]{\cd{/}} (con.west);
\draw[flecha] (ngx.east) -- node[etiq,pos=0.68,above]{\cd{/predictivo/}} (pre.west);
\draw[flecha] (ngx.east) -- node[etiq,pos=0.72,below,sloped]{\cd{/historico/}} (his.west);

\draw[flechad] (con.south) -- node[etiq,right]{loopback} (pre.north);
\draw[flecha] (pg.east) -- ++(6mm,0) |- node[etiq,pos=0.25,right]{fuente} (pre.east);
\draw[flecha] (tim.south) -- ++(0,-5mm) -| node[etiq,pos=0.25,below]{\cd{docker exec}} (pre.south);

\draw[flecha] (pre.east) -- ++(8mm,0) |- (sb.north);
\draw[flecha] (his.east) -- (sb.west);

\end{tikzpicture}
}
\caption{Qué corre dónde. Los dos agentes y la réplica escuchan sólo en el
\emph{loopback}; lo único publicado es nginx, con TLS, y detrás de él la consola con su
control de acceso y las dos rutas de servicio que consume VisioneFlow.}
\end{figure}

\subsection{Por qué la plataforma vive en su propio servidor}

Hasta el 2026-08-24 los dos microservicios corrían \emph{dentro} de la EC2 de VisioneFlow,
la plataforma de agentes de la empresa. Funcionaba, pero ataba dos proyectos que no
tienen nada que ver: el despliegue del pronóstico vivía en el repositorio de la
plataforma, cada despliegue del \emph{backend} hacía \cd{git reset --hard} y había que
acordarse de no tocar ciertos archivos, y una caída de cualquiera de los dos arrastraba
al otro.

Desde entonces la plataforma de datos tiene servidor propio, dominio propio y
certificado propio. VisioneFlow le sigue hablando, pero ahora como lo que es: una API
externa, con su clave, sobre HTTPS. Es la misma relación que tendría con cualquier
proveedor de terceros, y esa es exactamente la idea.

\subsection{Qué corre dónde}

\begin{center}
\small
\begin{tabular}{@{}p{34mm}p{50mm}p{56mm}@{}}
\toprule
\textbf{Pieza} & \textbf{Dónde} & \textbf{Cómo sobrevive a un reinicio} \\
\midrule
Agente Predictivo & \cd{predictivo-predictivo-1} \newline \nota{\cd{127.0.0.1:8000}} & \cd{restart: unless-stopped} \\
Agente Histórico & \cd{historico-historico-1} \newline \nota{\cd{127.0.0.1:8010}} & ídem \\
Réplica de AgroDash & \cd{agrodash-pg} \newline \nota{\cd{127.0.0.1:5433}} & ídem, sobre volumen persistente \\
Consola & proceso Node \newline \nota{\cd{127.0.0.1:3001}} & gestor de procesos con estado guardado \\
Borde HTTPS & nginx del sistema \newline \nota{puertos 80 y 443} & unidad systemd habilitada \\
Ingesta cada 15 min & \cd{predictivo-etl.timer} & unidad systemd habilitada \\
Refresco cada 6 h & \cd{predictivo-refresh.timer} & ídem \\
Almacenamiento & Supabase & servicio gestionado \\
\bottomrule
\end{tabular}
\end{center}

\begin{aviso}[title=El servidor no está encendido las 24 horas]
Un programador de tareas de AWS apaga la instancia a las 19:00 y la enciende a las 07:00,
de lunes a viernes. De viernes a lunes está apagada el fin de semana completo. Es una
decisión de costo, no un fallo, pero condiciona todo lo demás: cualquier consumidor
externo tiene que disparar dentro de esa ventana. La ingesta no pierde nada porque sus
temporizadores llevan \cd{Persistent=true} y \cd{OnBootSec}, así que al encender se pone
al día sola con lo que se perdió.
\end{aviso}

\subsection{Qué se publica, y qué no}

Sólo hay tres cosas alcanzables desde internet, todas por el mismo nginx con TLS:

\begin{center}
\small
\begin{tabular}{@{}p{40mm}p{40mm}p{58mm}@{}}
\toprule
\textbf{Ruta} & \textbf{Va a} & \textbf{Quién puede} \\
\midrule
\cd{/} & la consola & Quien tenga la contraseña. Sin sesión válida, las
  páginas redirigen y \cd{/api/*} responde 401. \\
\cd{/predictivo/...} & el Predictivo & Quien tenga la clave de servicio. No pasa por el
  control de la consola: un flujo automático no tiene sesión. \\
\cd{/historico/...} & el Histórico & ídem. \\
\bottomrule
\end{tabular}
\end{center}

Los nombres anteriores (\cd{/forecast/} y \cd{/analizador/}) siguen respondiendo como
alias, para que el cambio de servidor no obligara a reconfigurar rutas en el canvas de
la plataforma.

\subsection{Por qué hay una réplica de AgroDash}

El servidor de Cartago está caído, así que la fuente de la ingesta pasó a ser una
réplica restaurada de un volcado, montada en el propio servidor. El puerto queda atado al
loopback y ocupa unos 6\,GB. Volver a la fuente viva es cambiar una línea de
configuración: la URL de Cartago quedó comentada, no borrada.

Un detalle que cuesta tiempo si no se sabe: si el \cd{pg\_restore} de la máquina es más
nuevo que el servidor, falla con \cd{unrecognized configuration parameter "transaction\_timeout"}.
Por eso la restauración se hace con el \cd{pg\_restore} de \emph{adentro} del contenedor.

Al mudar el servidor, la réplica se movió por la red privada de la nube, sin salir a
internet, y se verificó contando las filas de las dos puntas: 21\,314\,662 lecturas en el
origen y las mismas 21\,314\,662 en el destino, con idéntica marca de tiempo máxima. Un
\cd{pg\_restore} que termina con código cero no es prueba de nada; el conteo sí.

\subsection{Salud del sistema}

Se mira sin entrar al servidor, por HTTP: \cd{/salud/ingesta} da la antigüedad de los
datos por variable (503 si están viejos) y \cd{/salud/panel} agrega errores recientes,
gasto del día y última predicción. La consola los muestra en su sección de salud.

\begin{aviso}[title=El 503 actual es correcto]
La fuente de San Carlos está congelada desde el 2026-07-23: el subsistema de cajas con
sufijo \cd{SC} y los nodos de suelo dejaron de reportar. \cd{/salud/ingesta} responde 503
porque está informando la antigüedad \emph{real} de los datos, no porque el agente falle.
Todo corre en modo ``sólo histórico''; cuando el equipo restaure la ingesta, el ETL
rellena solo por ser idempotente e incremental.
\end{aviso}

\subsection{Lección de observabilidad}

El ETL falló cada 15 minutos durante nueve días \textbf{sin dejar rastro} en
\cd{agente\_log}. La causa: la conexión a la fuente estaba \emph{fuera} del bloque
\cd{try}, así que la excepción se propagaba antes de que hubiera con qué registrarla.
La corrección tiene tres partes, y las tres importan: conectar primero al store (que es
donde se escribe el log), meter la conexión a la fuente dentro del \cd{try} con su propio
evento de fallo, y \textbf{volver a lanzar} la excepción para que systemd siga marcando
la unidad como fallida. Hay una prueba de regresión que fija ese comportamiento.

\subsection{Diagnóstico rápido}

\begin{center}
\small
\begin{tabular}{@{}p{56mm}p{84mm}@{}}
\toprule
\textbf{Síntoma} & \textbf{Dónde mirar} \\
\midrule
La ingesta no trae datos & \cd{systemctl status predictivo-etl.service} y la tabla
  \cd{agente\_log} filtrando por nivel de error. \\
El pronóstico responde datos viejos & \cd{GET /salud/ingesta}: la fuente está congelada,
  no es un fallo del agente. \\
Nada responde, de noche o en fin de semana & El servidor está apagado por programación.
  No es una caída. \\
Todo responde 429 & Tope diario de gasto, límite de ritmo del servicio, o el freno de
  nginx si es el acceso a la consola. \\
La consola responde 503 & Falta \cd{DEBUGGER\_PASSWORD} en el entorno. Falla cerrada a
  propósito. \\
Un servicio da 502 & El contenedor está caído; nginx sigue en pie. \\
Supabase rechaza escrituras & El plan gratuito son 500\,MB y está cerca del límite. \\
\bottomrule
\end{tabular}
\end{center}

\subsection{Pruebas}

Ninguna suite necesita credenciales ni red hacia Supabase; si una empieza a pedirlas,
dejó de ser unitaria. Corren en CI en cada \emph{push} y cada PR.

\begin{center}
\small
\begin{tabular}{@{}p{40mm}p{26mm}p{68mm}@{}}
\toprule
\textbf{Componente} & \textbf{Pruebas} & \textbf{Qué cubren} \\
\midrule
Agente Predictivo & 253 & Horizonte, física, herramienta, ETL, límites, salud, claves,
  riesgo, climatología, observabilidad, API. \\
Agente Histórico & 61 & API, costos, lista blanca de datos, control de calidad, cielo,
  contexto. \\
Consola & 3 guiones de humo & Construcción, control de acceso y formato de respuestas. \\
ETL de CSV & ninguna & Sin suite; la verificación es el \emph{notebook} de EDA. \\
\bottomrule
\end{tabular}
\end{center}
